% ============================================================
%  ch6_conclusion_future.tex
%  Chapter 6: Conclusions and Future Work
%  Master's Thesis -- SQL MATCH_RECOGNIZE on Pandas DataFrames
% ============================================================
\chapter{Conclusions and Future Work}
\label{chap:conclusions}

% ----------------------------------------------------------------
\section{Discussion}
\label{sec:discussion}
% ----------------------------------------------------------------

The evaluation shows that extending \texttt{MATCH\_RECOGNIZE} semantics
to \texttt{pandas} not only preserves SQL:2016 correctness but also
reduces developer effort for complex row-pattern recognition tasks.
To ensure correctness, the implementation was validated against Trino
and Oracle, confirming consistent results.  In contrast, a baseline
implementation in idiomatic \texttt{pandas} required hand-written
logic using \texttt{groupby}, \texttt{shift}/\texttt{diff}, and index
bookkeeping to emulate skip and overlap policies, underscoring the
additional complexity of reproducing the same semantics without native
support.  These comparisons emphasise the practical benefits of a
declarative approach: pattern logic can be expressed once in SQL and
executed directly on in-memory DataFrames, bridging SQL--Python
workflows.  Beyond correctness, the declarative integration also
promotes readability, maintainability, and interoperability, since
analysts can share queries across systems without rewriting custom
logic.  This lowers the barrier for adopting advanced row-pattern
recognition in data science pipelines, supports collaborative
workflows between SQL practitioners and Python developers, and
demonstrates how SQL-based semantics can be embedded into modern
analytics ecosystems.

From an architectural standpoint, the pipeline---ANTLR4-based parser,
\texttt{NFABuilder}, \texttt{DFABuilder}, and
\texttt{EnhancedMatcher}---is kept modular, which contributes to
semantic transparency and backend portability.  The engine already
incorporates a smart multi-level pattern cache (LRU-based, up to
10{,}000 entries) that avoids recompiling identical patterns, and a
parallel partition executor that processes independent
\texttt{PARTITION BY} groups concurrently using a thread pool.
This modularity also opens the door to future optimisation
opportunities such as guard pushdown (early elimination of impossible
transitions) and state pruning (discarding unreachable NFA states at
compile time), which could further reduce evaluation overhead for
complex patterns.  The same modular design would allow the matching
core to target other columnar engines (e.g., Polars, Apache Arrow)
with minimal changes.

Nevertheless, limitations remain.  Python's interpreter overhead,
particularly in high-ambiguity patterns, still constrains raw
throughput compared to native C++/Java SQL engines.  While
vectorisation and Cython/\texttt{numba} acceleration mitigate this,
the performance gap for highly nondeterministic workloads persists due
to potential NFA state explosion.  In addition, certain vendor-specific
features such as nested navigation functions or proprietary
optimisations are out of scope for this initial design.
Finally, although this work emphasises correctness and modularity,
scaling the approach into distributed execution frameworks
(e.g., Dask, Spark) will require additional development.

% ----------------------------------------------------------------
\section{Limitations}
\label{sec:conc-limitations}
% ----------------------------------------------------------------

Despite the system's comprehensive capabilities, several limitations
remain.

\textbf{Complex pattern and quantifier interactions.}
Although the system supports concatenation, alternation, grouping, and
standard quantifiers (\texttt{*}, \texttt{+}, \texttt{?},
\texttt{\{n,m\}}), certain combinations---particularly multiple greedy
quantifiers nested within groups (e.g., \texttt{(A+B*)+C?})---can
trigger exponential state-space growth during automata construction.
This issue primarily arises with three or more levels of nesting
combined with unbounded quantifiers; simpler patterns and bounded
quantifiers behave efficiently.

\textbf{Limited support for user-defined aggregate functions.}
While the engine supports a comprehensive set of built-in aggregates
(\texttt{SUM}, \texttt{COUNT}, \texttt{MIN}, \texttt{MAX},
\texttt{AVG}, \texttt{STDDEV}, \texttt{VARIANCE}, \texttt{ARRAY\_AGG},
\texttt{STRING\_AGG}, \texttt{MAX\_BY}/\texttt{MIN\_BY}, conditional
aggregates, and \texttt{RUNNING}/\texttt{FINAL} semantics), support
for fully arbitrary user-defined aggregate functions is limited.

\textbf{Performance on large datasets.}
The system performs efficiently on moderate-sized datasets but may
require additional optimisations for very large datasets.

\textbf{Memory usage for large patterns.}
Patterns with many variables and complex quantifiers can generate
large automata that increase memory consumption.

\textbf{Integration with query optimisers.}
Because the pattern-matching engine currently operates independently of
database query optimisers, it may miss plan-level optimisation
opportunities.

% ----------------------------------------------------------------
\section{Conclusion}
\label{sec:conclusion}
% ----------------------------------------------------------------

This thesis presented a SQL-in-\texttt{pandas} engine for executing
\texttt{MATCH\_RECOGNIZE} queries over DataFrames.  The engine accepts
SQL queries through an ANTLR4-generated parser (using an adapted Trino
grammar), compiles patterns to NFAs via \texttt{NFABuilder}, converts
them to DFAs via \texttt{DFABuilder}, and executes matches through
\texttt{EnhancedMatcher}---a thread-safe, backtracking-capable engine
that supports all four \texttt{AFTER MATCH SKIP} modes
(\texttt{PAST LAST ROW}, \texttt{TO NEXT ROW}, \texttt{TO FIRST},
\texttt{TO LAST}), both \texttt{ONE ROW PER MATCH} and
\texttt{ALL ROWS PER MATCH} output modes, navigation functions
(\texttt{FIRST}, \texttt{LAST}, \texttt{PREV}, \texttt{NEXT}),
\texttt{CLASSIFIER()}, \texttt{MATCH\_NUMBER()}, and a comprehensive
set of aggregate functions including \texttt{SUM}, \texttt{COUNT},
\texttt{MIN}, \texttt{MAX}, \texttt{AVG}, \texttt{STDDEV},
\texttt{ARRAY\_AGG}, conditional aggregates (\texttt{COUNT\_IF},
\texttt{SUM\_IF}), and \texttt{RUNNING}/\texttt{FINAL} semantics.
Results match vendor engines while substantially reducing developer effort.

\texttt{MATCH\_RECOGNIZE} allows data scientists and analysts to use
powerful pattern-matching semantics directly within their familiar
\texttt{pandas} environment, without the need for complex Python code or
external SQL engine dependencies.  This reduces development complexity
and enhances productivity for sequential data analysis across domains,
including financial analysis, log processing, and time series pattern
detection.

By enabling SQL-native pattern recognition directly in Python, the proposed
approach bridges the gap between the expressiveness of relational queries
and the flexibility of in-memory analytics.  This opens the door to
unified and portable pipelines that preserve both semantics and developer
productivity.

% ----------------------------------------------------------------
\section{Future Work}
\label{sec:conc-future}
% ----------------------------------------------------------------

Future work will extend the engine along three complementary directions.
First, \textbf{distributed execution} will be pursued by enabling the
matcher to scale across multi-node environments with boundary-aware
state propagation, ensuring efficiency on large datasets.  Second, a
\textbf{backend abstraction layer} will be developed to target
high-performance columnar engines such as Polars, Dask, Vaex, and
Apache Arrow, thereby further closing the performance gap with native
SQL systems.  Finally, \textbf{tooling and adoption} will be strengthened,
including developer utilities for pattern debugging, match
visualisation, and seamless SQL--Python interoperability within notebook
environments.  Together, these directions will enhance scalability,
performance, and usability, broadening the impact of the approach.

By addressing the identified limitations and implementing the planned
enhancements, the goal is to develop a more adaptable and efficient
solution capable of handling complex pattern-matching scenarios across
various data processing environments.  This roadmap provides a clear
path for improving performance and expanding the scope of the current
implementation.
